% =====================================================================
%  第 6 章 §6.4 同伦代数与形变理论（ch06b2）
%  页段：PDF 304–318（原书 p.294–308）
% =====================================================================

% [原文 p.294]
\section{同伦代数与形变理论}

\takeaway{把复结构与向量丛的形变方程 $\bar\partial_\epsilon\bar\partial_\epsilon=0$、$(\bar\partial+B)^2=0$ 提升为 $A_{\infty}$／$L_{\infty}$ 代数上的 Maurer–Cartan 方程 $m_*(e^b)=0$；形变空间 $\hat M(C)$ 及其模空间 $M(C)$ 在同伦等价下不变，而 $C$ 形式时小平–Spencer 写像满射。}

本节与下一节要给 B 模型情形下的观察 6.73 赋予数学内容。其基础是复结构的形变理论。从它的概要开始。详见 [209]、[210]、[371]。先考虑讨论稍简单的向量丛的复结构形变。设 $E\to M$ 是复流形 $M$ 上的向量丛，并设 $E$ 已给定全纯向量丛的结构。\footnote{复向量丛只是指纤维为复向量空间；全纯向量丛则指转换函数是全纯（holomorphic）的。}把作为出发点的这个结构记作 $\bar\partial\colon E\to E\otimes\Lambda^{0,1}M$。考察把 $\bar\partial$ 形变而得的 $E$ 上全纯向量丛的结构。设 $\bar\partial_\epsilon$ 是把 $\bar\partial$ 形变而得的全纯向量丛结构的族。由引理 2.51，$\bar\partial_\epsilon-\bar\partial$ 是 0 阶微分算子。于是展开为
\begin{equation}
  \bar\partial_\epsilon=\bar\partial+\epsilon B_1+\epsilon^2B_2+\cdots \tag{6.28}
\end{equation}
这里 $B_i$ 是向量丛 $\operatorname{Hom}(E,\Lambda^{0,1}\otimes E)$ 的截面。$\bar\partial_\epsilon$ 给出全纯向量丛结构的充要条件，已在定理 2.49 中说明，即可积条件 $\bar\partial_\epsilon\bar\partial_\epsilon=0$。代入 (6.28)，逐次看 $\epsilon$ 的各幂的系数，便得到下面一列方程。
\begin{equation}
  \bar\partial\bar\partial=0, \tag{6.29}
\end{equation}
\begin{equation}
  \bar\partial B_1=0, \tag{6.30}
\end{equation}
\begin{equation}
  \bar\partial B_2+B_1\wedge B_1=0, \tag{6.31}
\end{equation}
\begin{equation}
  \bar\partial B_3+B_1\wedge B_2+B_2\wedge B_1=0, \tag{6.32}
\end{equation}
$\cdots$

在考察这个方程组之前，先看清两个形变 $\bar\partial+B_\epsilon$ 与 $\bar\partial+B'_\epsilon$ 相同是指什么条件。

\begin{definition}{6.105}
  若存在 $\varphi_\epsilon\colon E\to E$ 作为复向量丛的自同构的族，使
% [原文 p.295]
  \begin{equation}
    \varphi_\epsilon\circ(\bar\partial+B_\epsilon)=(\bar\partial+B'_\epsilon)\circ\varphi_\epsilon \tag{6.33}
  \end{equation}
  成立，则称这两个形变是同构的。
\end{definition}

把 $\varphi_\epsilon=1+\epsilon g_1+\epsilon^2g_2+\cdots$ 作幂级数展开。$g_i$ 是 $\operatorname{Hom}(E,E)$ 的截面。代入 (6.33) 便得到 $\bar\partial(g_1)=B_1-B'_1$ 等。

回到方程组 (6.29)～(6.32)。(6.29) 是 $\bar\partial$ 的可积性，由假定已成立。由 (6.30) 可定出 Dolbeault 上同调 $[B_1]\in H^{1}_{\bar\partial}(M;\operatorname{Hom}(E,E))$ 的一个元素。由 $\bar\partial(g_1)=B_1-B'_1$ 可知，上同调类 $[B_1]$ 只由形变的同构类决定。也就是说，$H^{1}_{\bar\partial}(M;\operatorname{Hom}(E,E))$ 是方程 $\bar\partial_\epsilon\bar\partial_\epsilon=0$ 的线性化方程的解空间。下面的定理是已知的。

\begin{theorem}{6.106}
  若 $H^{2}_{\bar\partial}(M;\operatorname{Hom}(E,E))=0$，则存在以 $H^{1}_{\bar\partial}(M;\operatorname{Hom}(E,E))=0$ 的 $0$ 的近旁为参数的、$E$ 的全纯向量丛结构的族。
\end{theorem}

% [Tutor Note: 定理 6.106 中“$H^{1}_{\bar\partial}(\cdots)=0$ 的 0 的近旁”按上下文应为“$H^{1}_{\bar\partial}(\cdots)$ 的 0 的近旁”，即参数空间取在 $H^1$ 里而非取零空间；原书 p.295 此处多印了“$=0$”，疑为笔误，照录。]

下面考察 (6.31)。这个条件是积 $[B_1]\wedge[B_1]\in H^{2}_{\bar\partial}(M;\operatorname{Hom}(E,E))$ 为 0。回忆记号：(6.31) 的 $B_1\wedge B_1$ 在微分形式上取外积，在 $\operatorname{Hom}(E,E)$ 的元素上取合成。这个定义可推广为：给定三个全纯向量丛 $E_1,E_2,E_3$ 时的映射
\begin{align*}
  &H^{k_1}_{\bar\partial}(M;\operatorname{Hom}(E_1,E_2))\otimes H^{k_2}_{\bar\partial}(M;\operatorname{Hom}(E_2,E_3))\\
  &\qquad\longrightarrow H^{k_1+k_2}_{\bar\partial}(M;\operatorname{Hom}(E_1,E_3)).
\end{align*}
这称为\textbf{米田积}（Yoneda product）。\footnote{米田积可推广到 $E_i$ 未必是全纯向量丛、而只是解析凝聚层的情形。那时不用 $H^{k_i}_{\bar\partial}(M;\operatorname{Hom}(E_1,E_2))$，而考虑 $\operatorname{Ext}^{k_i}(E_1,E_2)$。}当 $E$ 是线丛时，$\operatorname{Hom}(E,E)$ 中的合成可换，故由外积的分次可换性得 $B_1\wedge B_1=-B_1\wedge B_1=0$。当 $E$ 的秩不小于 2 时，$\operatorname{Hom}(E,E)$ 非可换，于是 $[B_1]\wedge[B_1]$ 可能不为 0。

(6.31) 意味着这个积为 0。也就是说，映射 $H^{1}(M;\operatorname{Hom}(E,E))\to H^{2}(M;\operatorname{Hom}(E,E))$，$[B_1]\mapsto[B_1\wedge B_1]$ 给出形变的障碍。

% [原文 p.296]
为说明这个映射的意义，写下定理 6.106 的推广。

\begin{theorem}{6.107}
  存在在原点近旁定义的全纯映射 $s\colon H^{1}_{\bar\partial}(M;\operatorname{Hom}(E,E))\to H^{2}_{\bar\partial}(M;\operatorname{Hom}(E,E))$，且 $s^{-1}(0)$ 与 $E$ 的全纯结构的模空间中 $\bar\partial$ 的近旁（作为复解析空间或概形）同型。
\end{theorem}

称 $s$ 为\textbf{仓西写像}（Kuranishi map）。映射 $[B_1]\mapsto[B_1\wedge B_1]$ 是仓西写像的 2 次近似。（仓西写像在 0 处的值以及 0 处的 1 阶微分都是 0。）

考察方程 (6.32)。(6.32) 意味着 $[B_1\wedge B_2+B_2\wedge B_1]\in H^{2}(M;\operatorname{Hom}(E,E))$ 为 0。

一般地，设有 $[a_i]\in H^{k_i}(M;\operatorname{Hom}(E_i,E_{i+1}))$ 且 $[a_1\wedge a_2]=0$，$[a_2\wedge a_3]=0$。此时取 $b_1,b_2$ 使 $\bar\partial b_1=a_1\wedge a_2$，$\bar\partial b_2=a_2\wedge a_3$，考虑 $b_1\wedge a_3+(-1)^{k_1}a_1\wedge b_2$，则可验证 $\bar\partial(b_1\wedge a_3+(-1)^{k_1}a_1\wedge b_2)=0$，于是得到 $H^{k_1+k_2+k_3-1}(M;\operatorname{Hom}(E_1,E_3))$ 的元素 $[b_1\wedge a_3+(-1)^{k_1}a_1\wedge b_2]$。这称为\textbf{Massey–米田 3 重积}（Massey–Yoneda triple product）。不过，Massey–米田 3 重积随 $b_1,b_2$ 的取法而变：从 $H^{k_1-1}(M;\operatorname{Hom}(E_1,E_2))\oplus H^{k_2-1}(M;\operatorname{Hom}(E_2,E_3))$ 到 $H^{k_1+k_2+k_3-1}(M;\operatorname{Hom}(E_1,E_3))$ 的映射 $x_1\oplus x_2\mapsto x_1\wedge a_3+a_1\wedge x_2$ 的像的部分不能确定。

由以上所述，定理 6.107 的仓西写像的 3 次项由 Massey–米田 3 重积给出。4 次以后的项不再叙述，它们是高次 Massey–米田积。

以上讨论对 $M$ 的复结构（而不是其上的向量丛）的模空间也几乎可以平行进行。\footnote{复结构的形变情形在历史上更早。}
\[
  b=\sum b_{i,j}\frac{\partial}{\partial z_i}\otimes d\bar z_j\in\Gamma(M;TM^{1,0}\otimes\Lambda^{0,1}M)
\]
设 $\bar\partial\colon C^\infty(M)\to\Gamma(M;\Lambda^{0,1}M)$ 表示原来的复结构，用下式定义 $\bar\partial+b\colon C^\infty(M)\to\Gamma(M;\Lambda^{0,1}M)$：
\begin{equation}
  (\bar\partial+b)(f)=\bar\partial(f)+\sum b_{i,j}\frac{\partial f}{\partial z_i}\otimes d\bar z_j. \tag{6.34}
\end{equation}

% [原文 p.297]
\begin{remark}{6.108}
  在向量丛的形变情形，含 $B$ 的项是 0 阶微分算子，即 $\bar\partial_\epsilon-\bar\partial$ 不含微分。(6.34) 中含 $b$ 的项是 1 阶微分算子。这是很大的差别。
\end{remark}

来计算 $\bar\partial+b$ 确定复结构这一条件式 $(\bar\partial+b)(\bar\partial+b)=0$。

\begin{definition}{6.109}
  用下式定义 $[b,b']\in\Gamma(M;TM^{1,0}\otimes\Lambda^{0,2}M)$：
  \[
    [b,b']=\sum\Bigl(b_{i,j}\frac{\partial b'_{i',j'}}{\partial z_i}
    -b'_{i',j'}\frac{\partial b_{i,j}}{\partial z_i}\Bigr)\otimes d\bar z_j\wedge d\bar z_{j'}.
  \]
\end{definition}

\begin{lemma}{6.110}
  $(\bar\partial+b)(\bar\partial+b)=0$ 与 $\bar\partial b+\frac{1}{2}[b,b]=0$ 等价。
\end{lemma}

证明可由直接计算得出。参见 [371]。置 $b_\epsilon=\epsilon b_1+\epsilon^2b_2+\cdots$ 计算 $\bar\partial b+\frac{1}{2}[b,b]=0$，便得到
\begin{equation}
  \begin{cases}
    \bar\partial\bar\partial=0,\\
    \bar\partial b_1=0,\\
    \bar\partial b_2+\tfrac{1}{2}[b_1,b_1]=0,\\
    \cdots
  \end{cases} \tag{6.35}
\end{equation}
% [Tutor Note: (6.35) 第二式带系数 $1/2$，而向量丛情形的 (6.31) 不带。原因是定义 6.109 的 $[b,b']$ 未作对称化，$1/2$ 正是为抵消把 $[b,b]$ 的两个变量各算一次所造成的重复；若改用对称化的括号，$1/2$ 便并入其中。]

由此可与向量丛的形变作同样的考察。与定理 6.106 相当的定理由 Kodaira–Nirenberg–Spencer [208] 给出，与定理 6.107 相当的定理由 Kuranishi [229] 给出。

在向量丛的形变情形出现了结合且一般非可换的代数（algebra）$\operatorname{Hom}(E,E)$。在复结构的形变情形，充当这一角色的是李代数（Lie algebra）$\Gamma(M;TM^{0,1})$。（定义 6.109 是向量场的括号积与微分形式的外积的组合。）

本节的主要目的是用同调代数（同伦代数）的语言讨论形变理论。在抽象的同调代数框架中讨论，可以把 §6.5 所论的两个情形——B 模型与形变量子化——以及 §6.6 所论的拉格朗日子流形的模空间的量子形变这三者，用统一的框架

% [原文 p.298]
来处理。\footnote{这三种类似并非偶然，作者认为它们会融合成一个大的理论，但目前还在中途。}

同伦代数式的形变理论有李代数版与代数版，分别对应复结构的形变与向量丛的形变。以下先把微分李代数、微分代数双方都推广为强同伦李代数（strong homotopy Lie algebra）又名 $L_{\infty}$ 代数（$L_{\infty}$ algebra），以及强同伦代数（$A_{\infty}$ 代数）。这样做有两个好处。其一是后述的定理 6.123：即使处理微分李代数、微分代数，映射也不只限于微分李代数的映射、微分代数的映射，而扩大为 $L_{\infty}$ 映射、$A_{\infty}$ 映射更为方便。其二是由于量子效应，微分代数或微分李代数变成 $A_{\infty}$ 代数或 $L_{\infty}$ 代数的例子是存在的。前者的例子是 §6.6 所论的拉格朗日子流形的 Floer 同调，后者的例子是（本书不处理）Eliashberg–Hofer 等人的接触同调（contact homology）（论文尚未发表，预告见 [94]）。

以下叙述参考了 Schechtman [324]\footnote{据说受了 Drinfeld 1988 年的私人通信与 Kontsevich 的讲演的启发。}、Schlessinger [325]、Schlessinger–Stasheff [326]、Goldman–Millson [152], [153]、Kontsevich [221]、Manin [243] 等。包含障碍类的处理，我想 [122] 是最早的。

设 $R$ 是 $\C$ 上的可换环，$C$ 是分次 $R$ 模，使用前一节的记号。用 $\Pi C$ 表示把 $C$ 的次数平移后的对象，即置 $(\Pi C)^k=C^{k+1}$。（因为在上同调中考虑，次数写在上方。）把定义 6.66 中引入的 $B\Pi C$、$E\Pi C$ 的张量积为 $k$ 个的部分写作 $B_k\Pi C$、$E_k\Pi C$。考虑次数为 $+1$ 的映射列
\[
  m_k\colon B_k\Pi C\to\Pi C,\qquad l_k\colon E_k\Pi C\to\Pi C.
\]

\begin{remark}{6.111}
  按平移前的原次数考虑，$m_k,l_k$ 是 $2-k$ 次的。可把 $m_1$ 看成余边界、$m_2$ 看成杯积之类的东西：$m_1$ 是 1 次、$m_2$ 是 0 次，故说法相符。特意平移次数，其目的之一是把符号系统地定下来。
\end{remark}

% [原文 p.299]
$B$ 的情形用 (6.36)，$E$ 的情形用 (6.37)，把 $B\Pi C\to B\Pi C$、$E\Pi C\to E\Pi C$ 分别扩张。
\begin{equation}
  \hat d_k(x_1\cdots x_n)=\sum_{i=1}^{n-k+1}\pm x_1\otimes\cdots\otimes m_k(x_i\cdots x_{i+k-1})\otimes\cdots\otimes x_n, \tag{6.36}
\end{equation}
\begin{equation}
  \hat d_k(x_1\cdots x_n)=\sum_k\sum_{\sigma\in\mathfrak{S}_n}\pm\frac{1}{k!(n-k)!}
  l_k(x_{\sigma(1)}\cdots x_{\sigma(k)})\otimes x_{\sigma(k+1)}\otimes\cdots\otimes x_{\sigma(n)}. \tag{6.37}
\end{equation}
这里符号在 (6.36) 中由 $(-1)^{\deg x_1+\cdots+\deg x_{k-1}}$（$\deg$ 为平移后的次数）决定，在 (6.37) 中由 $x_1\cdots x_n=\pm x_{\sigma(1)}\cdots x_{\sigma(n)}$ 决定。又 $\mathfrak{S}_n$ 表示 $\{1,\dots,n\}$ 之间的置换全体。$k$ 取 0 以上的整数与取 1 以上的整数这两种情形都考虑。置 $\hat d=\sum\hat d_k$。

$BC,EC$ 上有余积（coproduct）$\Delta\colon BC\to BC\otimes BC$，$\Delta\colon EC\to EC\otimes EC$。即
\begin{equation}
  \Delta(x_1\cdots x_n)=\sum_{k=0}^{n}(x_1\cdots x_k)\otimes(x_{k+1}\cdots x_n), \tag{6.38}
\end{equation}
\begin{equation}
  \Delta(x_1\cdots x_n)=\sum_{\sigma\in\mathfrak{S}_n}\sum_{k=0}^{n}\pm\frac{1}{k!(n-k)!}
  (x_{\sigma(1)}\cdots x_{\sigma(k)})\otimes(x_{\sigma(k+1)}\cdots x_{\sigma(n)}) \tag{6.39}
\end{equation}
（(6.39) 的符号与 (6.37) 相同）。（$k=0$ 时把 $x_1\cdots x_k$ 看作 1。）

% [Tutor Note: (6.37)、(6.39) 中的 $\pm$ 统一由 $x_1\cdots x_n=\pm x_{\sigma(1)}\cdots x_{\sigma(n)}$ 确定，而因子 $1/(k!(n-k)!)$ 抵消了同一分法被数 $k!(n-k)!$ 遍的重数。$B$ 的情形（(6.36)、(6.38)）没有 $\pm$，因为那里的张量积是有序的，不需要对称化。]

$\Delta$ 是余结合的（coassociative），即 $(\Delta\otimes1)\Delta=(1\otimes\Delta)\Delta$。又 $(EC,\Delta)$ 是余可换的（cocommutative），即定义 $T\colon EC\otimes EC\to EC\otimes EC$ 为 $T(x\otimes y)=(-1)^{xy}y\otimes x$ 时，有 $T\Delta=\Delta$。（今后，$(-1)$ 的肩上出现 $x$ 时，$x$ 指 $\deg x$。）$BC$、$EC$ 分别是自由余结合余代数（free coassociative coalgebra）与自由余可换余结合余代数（free coassociative cocommutative coalgebra）。

又 $\hat d$ 是余微分（coderivative），即 $\Delta\hat d=(\hat d\otimes1+1\otimes\hat d)\Delta$

% [原文 p.300]
成立。（置 $(1\otimes\hat d)(x\otimes y)=(-1)^xx\otimes\hat d(y)$。）

\begin{definition}{6.112}
  称 $(B\Pi C,m_k)$（$k>0$）为 \textbf{$A_{\infty}$ 代数}，是指 $\hat d\hat d=0$。称 $(E\Pi C,l_k)$（$k>0$）为 \textbf{$L_{\infty}$ 代数}，也是指 $\hat d\hat d=0$。$k\ge0$ 的情形称为\textbf{弱 $A_{\infty}$ 代数}（weak $A_{\infty}$ algebra）、\textbf{弱 $L_{\infty}$ 代数}（weak $L_{\infty}$ algebra）。\footnote{这一带的术语尚无标准，文献中 $k\ge0$ 与 $k>0$ 两种都有。}
\end{definition}

\begin{remark}{6.113}
  $A_{\infty}$ 代数的概念由 Stasheff [347] 引入。
\end{remark}

来看 $m_0=0$ 的情形下 $\hat d\hat d=0$ 的具体关系式。假定 $m_0=0$，则 $B\Pi C$、$E\Pi C$ 的滤过就确定了。\footnote{这个滤过与下面要讲的 Novikov 环的能量滤过不同，须注意。}即置
\begin{equation}
  \mathfrak{G}_kB\Pi C=\sum_{\ell\le k}B_\ell\Pi C,\qquad
  \mathfrak{G}_kE\Pi C=\sum_{\ell\le k}E_\ell\Pi C. \tag{6.40}
\end{equation}
称这个滤过为\textbf{个数滤过}（number filter）。此时 $\hat d(\mathfrak{G}_kB\Pi C)\subset\mathfrak{G}_kB\Pi C$，$\hat d(\mathfrak{G}_kE\Pi C)\subset\mathfrak{G}_kE\Pi C$ 成立。（若 $m_0,l_0$ 不为 0，则 $\hat d$ 不保持个数滤过。）在 $\mathfrak{G}_kB\Pi C$ 上考察 $\hat d^2=0$，就依次得到关系式，即
\begin{equation}
  0=m_1m_1, \tag{6.41}
\end{equation}
\begin{equation}
  0=m_1m_2(xy)+m_2(m_1(x)y)+(-1)^xm_2(xm_1(y)), \tag{6.42}
\end{equation}
\begin{equation}
  \begin{aligned}
    0={}&m_1m_3(xyz)+m_3(m_1(x)yz)\\
    &+(-1)^xm_3(xm_1(y)z)+(-1)^{x+y}m_3(xym_1(z))\\
    &+m_2(m_2(xy)z)+(-1)^xm_2(xm_2(yz))
  \end{aligned} \tag{6.43}
\end{equation}
成立。(6.41) 表示 $m_1$ 是微分，(6.42) 表示积 $m_2$ 满足积的微分法公式。(6.43) 是结合律，但符号须注意。即置 $x\cdot y=(-1)^{x-1}m_2(xy)$。\footnote{据 Getzler–Jones [140]。}（$\deg x-1$ 是 $x$ 原本

% [原文 p.301]
的次数。）于是，例如 $m_3=0$ 时，(6.43) 就成为结合律 $(x\cdot y)\cdot z=x\cdot(y\cdot z)$。

同样计算可知，$L_{\infty}$ 代数的关系式中第 3 个是 Jacobi 律。即设 $l_2$ 以外全为 0：
\[
  l_2(xy)=(-1)^{xy}l_2(yx),
\]
\[
  l_2(l_2(xy),z)+(-1)^{(x+y)z}l_2(l_2(zx),y)+(-1)^{(y+z)x}l_2(l_2(yz),x)=0
\]
是 $L_{\infty}$ 关系式，而置 $[x,y]=(-1)^{x-1}l_2(xy)$ 后，它成为
\[
  [x,y]+(-1)^{(x+1)(y+1)}[y,x]=0,
\]
\[
  [[x,y],z]+(-1)^{(x+y)(z+1)}[[z,x],y]+(-1)^{(y+z)(x+1)}[[y,z],x]=0.
\]
这是关于平移前的原次数的分次李代数关系式。

由个数滤过得到下列谱序列。
\begin{equation}
  E_2\simeq BH^*(\Pi C,m_1)\Longrightarrow H^*(B\Pi C,\hat d), \tag{6.44}
\end{equation}
\begin{equation}
  E_2\simeq EH^*(\Pi C,m_1)\Longrightarrow H^*(E\Pi C,\hat d). \tag{6.45}
\end{equation}

% [Tutor Note: (6.44)、(6.45) 的 $E_2$ 项由 $m_1$ 在 $\Pi C$ 上的同调按张量积组装而成；高阶 $m_k$（$k\ge2$）正是谱序列的微分来源。因此“$C$ 是形式的”等价于说这些微分全为零，谱序列在 $E_2$ 退化——这正是定理 6.136 与定理 6.107 的分水岭。]

$L_{\infty}$ 代数、$A_{\infty}$ 代数的形式形变理论（即幂级数范围内的形变理论）中，不用 $R$，而用 Novikov 环或其部分环——幂级数环 $\C[[T]][u,u^{-1}]$。设 $\Lambda$ 为这两者之一。设 $\Lambda_{\mathrm{nov},0}$ 为 Novikov 环中由 $\sum T^{\lambda_i}\otimes u^{n_i}$，$\lambda_i>0$ 者全体构成的理想。置 $\Lambda_0=\Lambda\cap\Lambda_{\mathrm{nov},0}$。$\Lambda/\Lambda_0$ 与 $\C[u,u^{-1}]$ 同型。

$\Lambda$ 模 $C$ 的\textbf{赋值}（valuation）是指 $|\cdot|\colon C\to\R$，满足 $|ax|\ge|a|+|x|$ 对 $a\in\Lambda$，$x\in C$ 成立。这里 $|a|$ 是引理 6.77 之前定义的\textbf{能量}。由 $|\cdot|$ 给出的滤过称为\textbf{能量滤过}（energy filter）。（$\Lambda$ 为幂级数环时，设 $|\cdot|$ 取整数值。）

给定分次 $R=\C$ 模 $\overline C$ 与 $L_{\infty}$ 代数、$A_{\infty}$ 代数 $(E\Pi\overline C,\overline m_k)$、$(B\Pi\overline C,\overline m_k)$，且
\[
  \frac{C}{\Lambda_0C}\simeq\overline C\otimes\C[u,u^{-1}]
\]

% [原文 p.302]
者。以下假定 $C$ 由 $|x|=0$ 的元素 $x$ 在 $\Lambda_{\mathrm{nov}}$ 上生成。设 $m_k\colon\widehat{B\Pi C}\to C$ 或 $m_k\colon\widehat{E\Pi C}\to C$ 已给定，且 $m_k$ 保持能量滤过。此时完备化后考虑 $\hat d\colon\widehat{B\Pi C}\to\widehat{B\Pi C}$，$\hat d\colon\widehat{E\Pi C}\to\widehat{E\Pi C}$，而式 (6.36)、(6.37) 不变。由于 $m_k$ 保持滤过，$\hat d$ 可扩张到完备化上。

\begin{definition}{6.114}
  称 $(\widehat{B\Pi C},m_k)$ 为 $(\overline{B\Pi C},\overline m_k)$ 的形变、且是\textbf{滤过 $A_{\infty}$ 代数}（filtered $A_{\infty}$ algebra），是指 $\hat d\hat d=0$ 且把 $m_k$ 模 $\Lambda_{\mathrm{nov},0}$ 考虑时，它与把 $\overline m_k$ 系数扩张到 $\C[u,u^{-1}]$ 所得者一致。\textbf{滤过 $L_{\infty}$ 代数}（filtered $L_{\infty}$ algebra）也同样定义。把 $m_0,l_0$ 也包含在内考虑的情形，称为\textbf{滤过弱 $A_{\infty}$ 代数}（filtered weak $A_{\infty}$ algebra）、\textbf{滤过弱 $L_{\infty}$ 代数}（filtered weak $L_{\infty}$ algebra）。
\end{definition}

\begin{remark}{6.115}
  把 $A_{\infty}$ 代数、$L_{\infty}$ 代数用于形变理论，大致始于 [326]。滤过 $A_{\infty}$ 代数由 [122] 引入。
\end{remark}

设 $x_i$ 是 $C$ 的生成元且 $|x_i|=0$。$x_i$ 也可与 $\overline C$ 在 $\C$ 上的基等同看待。可表为
\begin{equation}
  m_k(x_{i_1}\cdots x_{i_k})=\sum_jm_{k,\lambda_j}(x_{i_1}\cdots x_{i_k})T^{\lambda_j}\otimes u^{n_j} \tag{6.46}
\end{equation}
（$l_{k,\lambda_j}$ 也同样定义）。$u$ 不出现的情形即 $n_j=0$ 的情形，若把数代入 $T$ 所得的 $m_k$ 收敛，就得到 $\overline C$ 上的 $A_{\infty}$ 代数、$L_{\infty}$ 代数结构。由假定，它在 $T=0$ 处与形变前的 $A_{\infty}$ 代数、$L_{\infty}$ 代数一致。

给定滤过弱 $A_{\infty}$ 代数、弱 $L_{\infty}$ 代数后，由此可构成滤过 $A_{\infty}$ 代数、$L_{\infty}$ 代数的模空间。说明这一构成。对 $b\in(\Pi C)^0=C^1$，$b\equiv0\ \mathrm{mod}\ \Lambda_0$，置
\begin{equation}
  e^b=1+b+b\otimes b+b\otimes b\otimes b+\cdots\qquad (A_{\infty}\ \text{代数情形}) \tag{6.47}
\end{equation}
\begin{equation}
  e^b=1+b+\frac{1}{2}b\otimes b+\frac{1}{3!}b\otimes b\otimes b+\cdots\qquad (L_{\infty}\ \text{代数情形}) \tag{6.48}
\end{equation}

% [原文 p.303]
由 $b\equiv0\ \mathrm{mod}\ \Lambda_0$ 与完备性，右端收敛。且 $\Delta e^b=e^b\otimes e^b$ 成立。

定义
\begin{equation}
  m_k^b(x_1\cdots x_k)=m_*(e^bx_1e^b\cdots e^bx_ke^b), \tag{6.49}
\end{equation}
\begin{equation}
  l_k^b(x_1\cdots x_k)=l_*(e^bx_1\cdots x_k) \tag{6.50}
\end{equation}
（$k\ge0$）。（在 (6.49)、(6.50) 中，把括号内的项展开，对每一项的 $*$ 代入使个数与括号中变量相合的 $k'\ge k$。）

% [Tutor Note: $m_*$、$l_*$ 指把全部 $m_k$、$l_k$ 合在一起、按张量积乘性延拓后的映射，故 $e^bx_1e^b\cdots$ 里的 $e^b$ 实际上是余积给出的若干张量因子。命题 6.116 说的就是“用 $b$ 扭一下”仍保持 $A_{\infty}$／$L_{\infty}$ 结构；这也是规范变换在代数一侧的实现。]

\begin{proposition}{6.116}
  $m_k^b$、$l_k^b$ 确定滤过弱 $A_{\infty}$ 代数或滤过弱 $L_{\infty}$ 代数的结构。
\end{proposition}

证明只需老实代入计算。留给读者（参见 [122]）。

\begin{proposition}{6.117}
  $m_k^b$ 成为滤过 $A_{\infty}$ 代数、$l_k^b$ 成为滤过 $L_{\infty}$ 代数的充要条件是 $m_*(e^b)=0$、$l_*(e^b)=0$。
\end{proposition}

证明由定义显然。

由 $\hat d(e^b)=e^bm_*(e^b)e^b$，命题 6.117 的条件也可写作 $\hat d(e^b)=0$。

把条件式 $m_*(e^b)=0$ 改写，例如 $A_{\infty}$ 代数情形就是
\begin{equation}
  m_0(1)+m_1(b)+m_2(b,b)+m_3(b,b,b)+\cdots=0. \tag{6.51}
\end{equation}
若原本 $m_0(1)=0$，则写作 $m_1=d$，$m_2=\wedge$ 时，(6.51) 成为 $db+b\wedge b+\cdots=0$。这在 $A_{\infty}$ 代数情形是 $(\bar\partial+B)^2=0$ 这一向量丛形变方程的类似物，在 $L_{\infty}$ 代数情形则是复结构形变方程的类似物。

(6.51)（及其 $L_{\infty}$ 代数情形的类似物）称为\textbf{莫雷尔–嘉当方程}（Maurer–Cartan equation），或称 \textbf{Batalin–Vilkovsky 主方程}（Batalin–Vilkovsky master equation）。

\begin{remark}{6.118}
  前者出自对称空间与李群的研究，后者出自约束系统的哈密顿力学。\footnote{可用日语读到的解说有 [273] 等，但那是物理学家所写，对以数学为专业的人不易读；数学家的文献有 [348] 等。}这些与第 3 章末尾只提过名字的 BRST 上同调等也有关。此外，外微分方程与哈密顿系统对称性方面的各种理论 [53] 以及 A. Vinogradov 的 $C$ 谱

% [原文 p.304]
序列 [377] 等也与此关系很深。
\end{remark}

\begin{definition}{6.119}
  用 $\hat M(C)$ 表示使 $\hat de^b=0$ 的 $b$ 全体。
\end{definition}

其次定义 $A_{\infty}$ 及 $L_{\infty}$ 代数的同伦同值。设有 $R$ 同态的族 $\varphi_k\colon B_k\Pi C\to\Pi C'$（或 $\varphi_k\colon E_k\Pi C\to\Pi C'$）。

\begin{lemma*}{6.120}
  存在唯一的余同态 $\hat\varphi\colon B\Pi C\to B\Pi C'$，使 $\hat\varphi$ 在 $B_k\Pi C$ 上的限制与投影 $B\Pi C'\to\Pi C'$ 的合成等于 $\varphi_k$。对 $\hat\varphi\colon E\Pi C\to E\Pi C'$ 也同样。
\end{lemma*}

\begin{proof}
  置
  \[
    \hat\varphi(x_1\cdots x_n)=\sum_{0=k_1\le\cdots\le k_\ell=n}
    \varphi_{k_2-k_1}(x_1\cdots x_{k_2})\cdots\varphi_{k_\ell-k_{\ell-1}}(x_{k_{\ell-1}+1}\cdots x_n)
  \]
  即可。它是余同态，即 $\Delta\hat\varphi=(\hat\varphi\otimes\hat\varphi)\Delta$ 容易看出。唯一性重看余同态的定义便知。$E$ 也同样。
\end{proof}

注意（$\varphi_0$ 不存在时）$\hat\varphi$ 保持个数滤过。

\begin{definition}{6.121}
  称 $\varphi_k$（$k\ge1$）为 \textbf{$A_{\infty}$ 同型写像}（$A_{\infty}$ homomorphism），是指 $\hat\varphi\hat d=\hat d\hat\varphi$。\textbf{$L_{\infty}$ 同型写像}（$L_{\infty}$ homomorphism）也同样定义。在 $k\ge0$ 的范围内考虑时，称为\textbf{弱 $A_{\infty}$ 同型写像}（weak $A_{\infty}$ homomorphism）、\textbf{弱 $L_{\infty}$ 同型写像}（weak $L_{\infty}$ homomorphism）。
\end{definition}

\begin{definition}{6.122}
  称 $A_{\infty}$ 写像是同伦同值写像，是指它在 $\hat d$ 的同调上诱导同构。$L_{\infty}$ 写像的情形也同样。
\end{definition}

“在同调上诱导同构”听起来或许比同伦同值的定义要弱，但拓扑的情形有 Whitehead 定理，即“诱导同调同构的（单连通流形之间的）连续映射是同伦同值”。我们的情形也有下面的定理。

\begin{theorem}{6.123}
  对 $A_{\infty}$ 写像 $\hat\varphi\colon B\Pi C\to B\Pi C'$，若 $\varphi_1\colon(\Pi C,m_1)\to(\Pi C',m'_1)$ 在同调上诱导同构，则存在 $A_{\infty}$ 写像 $\hat\varphi'\colon B\Pi C'\to B\Pi C$，且 $\hat\varphi'\hat\varphi$ 在同调上给出逆映射。$L_{\infty}$ 写像的情形也同样。
\end{theorem}

\begin{corollary}{6.124}
  $A_{\infty}$、$L_{\infty}$ 代数的同伦同值是等价关系。
\end{corollary}

% [原文 p.305]
\begin{remark}{6.125}
  定理 6.123 是考虑 $A_{\infty}$、$L_{\infty}$ 结构的重要好处。即使只在微分代数或微分李代数的范围内考虑，要使定理 6.123 成立，也必须一直考虑到 $A_{\infty}$、$L_{\infty}$ 写像。
  为构造导出范畴（derived category），要对范畴作局部化（localization），把所有在同调上诱导同构的映射变为可逆（参见 [136]、[173] 等）。若考虑 $A_{\infty}$、$L_{\infty}$ 写像，则这一操作不必要。
\end{remark}

\begin{remark}{6.126}
  由谱序列 (6.44)、(6.45)，若 $\varphi_1$ 在同调上诱导同构，则 $\hat\varphi$ 也在同调上诱导同构。进一步，$\hat\varphi$ 在 $\mathfrak{G}_kC$ 上的限制也在同调上诱导同构。
\end{remark}

\begin{proof}[定理 6.123 的证明]
  只证 $A_{\infty}$ 代数情形，$L_{\infty}$ 代数情形几乎相同。用关于 $k$ 的归纳法构成 $\varphi'_k$。对诱导链复形同构的链映射 $\varphi_1$，构成在同调上给出逆的 $\varphi'_1$ 是容易的。

  设已构成到 $\varphi'_{k-1}$。由 $\varphi'_1,\dots,\varphi'_{k-1}$ 定出 $\hat\varphi'_{k-1}\colon\mathfrak{G}_{k-1}B\Pi C'\to\mathfrak{G}_{k-1}B\Pi C$ 以及 $\bar\varphi'\colon\mathfrak{G}_kB\Pi C'/\mathfrak{G}_1B\Pi C'\to\mathfrak{G}_kB\Pi C/\mathfrak{G}_1B\Pi C$。由注意 6.126，$\hat\varphi'_{k-1}$、$\bar\varphi'$ 都在同调上诱导同构。定义 $\varphi'^0\colon\mathfrak{G}_kB\Pi C'\to\mathfrak{G}_kB\Pi C$ 为：在 $\mathfrak{G}_{k-1}B\Pi C'$ 上取 $\bar\varphi'$，在 $B_k\Pi C'$ 上取
  \[
    \varphi'^0(x_1\cdots x_k)=\sum_{\ell\ge2}\sum_{1=i_1<\cdots<i_\ell=k+1}
    \varphi_{i_2-i_1}(x_{i_1}\cdots x_{i_2-1})\cdots\varphi_{i_\ell-i_{\ell-1}}(x_{i_{\ell-1}}\cdots x_{i_\ell}).
  \]
  （$\varphi'_k$ 尚未定义，但一旦定义，就有 $\hat\varphi'=\varphi'_k+\varphi'^0$。）又置
  \[
    K=\{x\in B_k\Pi C\mid\hat dx\in\hat d(\mathfrak{G}_{k-1}\Pi C')\},\qquad
    B=\{m_1x\mid x\in B_k\Pi C\}.
  \]
  设 $B$ 的基为 $v_1,\dots,v_a$，$K$ 的基为 $v_1,\dots,v_a,v_{a+1},\dots,v_b$。令 $v_i=m_1w_i$，$i=1,\dots,a$，并置 $v_i+v'_i=\hat dw_i$，$i=1,\dots,a$。则 $v_1,\dots,v_b,w_1,\dots,w_a$ 成为 $B_k\Pi C$ 的基。对 $i=a+1,\dots,b$，把 $v'_i$ 取定使 $\hat d(v_i+v'_i)=0$。置 $\hat v_i=v_i+v'_i$（$i=1,\dots,b$）。

  由 $\bar\varphi\hat d=\hat d\bar\varphi$、$\varphi'^0\equiv\bar\varphi\ \mathrm{mod}\ \mathfrak{G}_1B\Pi C$ 可知 $\hat d\varphi'^0(\hat v_i)\in\mathfrak{G}_1B\Pi C=\Pi C$。故 $m_1\hat d\varphi'^0(\hat v_i)=\hat d\hat d\varphi'^0(\hat v_i)=0$ 成立。更强地，下列事实成立。

  \begin{lemma*}{6.127}
    $\hat d\varphi'^0(\hat v_i)\in\operatorname{Im}m_1$。
  \end{lemma*}

% [原文 p.306]
  \begin{proof}
    由于 $\bar\varphi\bar\varphi'$ 在同调上诱导恒等映射，存在 $z_i\in\Pi C$ 使 $\hat\varphi\hat\varphi'^0(\hat v_i)-\hat v_i-z_i\in\operatorname{Im}\hat d$。故 $\hat\varphi\hat d\hat\varphi'^0(\hat v_i)=\hat d\hat\varphi\hat\varphi'^0(\hat v_i)=m_1\hat z_i$。由于 $\hat\varphi$ 在 $m_1$ 的同调上诱导同构，得到本引理。
  \end{proof}

  由引理 6.127，存在 $y_i\in\Pi C$ 使
  \begin{equation}
    \hat d\varphi'^0(v_i+v'_i)+m_1y_i=0. \tag{6.52}
  \end{equation}
% [Tutor Note: 这段归纳法的结构值得记牢：先用 (6.52) 在 $\mathfrak{G}_{k-1}B\Pi C+K$ 上定出 $\varphi'_k$，再借它把 $\hat\varphi'$ 扩张到整个 $\mathfrak{G}_kB\Pi C$。两次定法之间相差 $\operatorname{Ker}m_1$ 的元素，这个自由度恰好在下一步（定义 $\varphi'_k(w_i)$ 时）用来消去障碍。]

  置 $\varphi'_k(v_i)=y_i$。于是 $\hat\varphi'$ 扩张到 $\mathfrak{G}_{k-1}B\Pi C+K$ 上，且在那里 $\hat d\hat\varphi'=\hat\varphi'\hat d$ 成立。$y_i$ 的取法有 $\operatorname{Ker}m_1$ 那么多的自由度。这个自由度在下一步用。

  定义 $\varphi'_k(w_i)$。由 $\hat d\bar\varphi=\bar\varphi\hat d$ 及 $\hat dw_i\equiv v_i\ \mathrm{mod}\ \mathfrak{G}_{k-1}B\Pi C$，
  \[
    \hat d\varphi'^0(w_i)-y_i-\varphi'^0(\hat dw_i)=\hat d\varphi'^0(w_i)-\hat\varphi'\hat d(w_i)\in\Pi C.
  \]
  另一方面，因 $\hat d(w_i)\in\mathfrak{G}_{k-1}B\Pi C+K$，有 $\hat d\bigl(\hat d\varphi'^0(w_i)-\hat\varphi'\hat d(w_i)\bigr)=0$。由此得 $m_1\bigl(\hat d\varphi'^0(w_i)-\hat\varphi'\hat d(w_i)\bigr)=0$。故把 $y_i$ 差一个 $\operatorname{Ker}m_1$ 的元素重新取定，可设 $\hat d\varphi'^0(w_i)-\hat\varphi'\hat d(w_i)=0$。于是置 $\varphi'_k(w_i)=0$ 而把 $\hat\varphi'$ 扩张到 $\mathfrak{G}_kB\Pi C$ 上，就有 $\hat d\hat\varphi'=\hat\varphi'\hat d$ 成立。至此归纳法完成。
\end{proof}

对 $A_{\infty}$、$L_{\infty}$ 结构在形变理论中的应用，重要的是形变模空间的同伦不变性。形变模空间如下定义。

\begin{definition}{6.128}
  设 $b,b'\in\hat M(C)$。称 $b$ 与 $b'$ \textbf{规范等价}（gauge equivalent），是指存在关于 $t$ 可微的 $b(t)\in\hat M(C)$ 与 $c(t)\in(\Pi C)_1$，满足
  \[
    \frac{\partial b(t)}{\partial t}=m_{1,b(t)}c(t)
  \]
  （$L_{\infty}$ 代数时以 $l_1$ 代替 $m_1$）。规范等价类全体用 $M(C)$ 表示。
\end{definition}

\begin{theorem}{6.129}
  若 $C$ 与 $C'$ 同伦同值，则 $M(C)$ 与 $M(C')$ 之间存在双射。
\end{theorem}

% [Tutor Note: 原书给“定理 6.129”加上“”引号，并在 p.306 自注“本来应在 $M(C)$ 上定出结构、再断言关于该结构同型，但作者至今不清楚应放入的正确结构（或可称为形式超概形）的精确定义，故以上述形式陈述定理”，说明这是作者自认尚未严格证明的陈述，与前后正式的定理区别对待。]

本来应当在 $M(C)$ 上定出结构、再断言关于该结构同型，但应放入的正确结构\footnote{或许可称为形式超概形（formal super scheme）。}的精确定义作者至今不明

% [原文 p.307]
了，故陈述了上述形式的定理。“定理 6.129”的证明可按下列步骤进行。

\textbf{方案 6.130}

(1) 对 $A_{\infty}$ 写像 $\hat\varphi\colon B\Pi C\to B\Pi C'$（或 $L_{\infty}$ 写像 $\hat\varphi\colon E\Pi C\to E\Pi C'$），构成写像 $\hat M(C)\to\hat M(C')$（定义 6.131），并证明它诱导写像 $M(C)\to M(C')$。

(2) 定义两个 $A_{\infty}$ 写像（$L_{\infty}$ 写像）同伦，并证明 $M(C)$ 上诱导的写像只依赖于同伦类。

(3) 对定理 6.123 的 $\hat\varphi'$，证明 $\hat\varphi'\hat\varphi$、$\hat\varphi\hat\varphi'$ 与恒等写像 $1_C$、$1_{C'}$ 同伦。这里 $1_C$ 指 $1_{C,1}\colon\Pi C\to\Pi C$ 为恒等写像、$1_{C,k}=0$（$k\ge2$）的 $A_{\infty}$、$L_{\infty}$ 写像。\hfill$\square$

利用定理 6.123，由方案 6.130 的 (1)、(2)、(3) 便得“定理 6.129”。实行方案 6.130 需要相当麻烦的同调代数。遗憾的是，本书放弃在一般情形下实行它，只给出 §6.5 所需的定理 6.136。§6.6 需要更一般的情形，那里所需的情形在 [122] 中有叙述。麻烦的是，把“定理 6.129”在 $A_{\infty}$、$L_{\infty}$ 代数的同调代数范畴中详细证明的文献似乎并不存在。概略在 [221] 中。不是 $L_{\infty}$ 代数、而在微分李代数范围内的证明见 [152]、[153] 等。[119]、[243]、[324]、[326] 也叙述了相关话题。

\begin{definition}{6.131}
  设 $\hat\varphi\colon B\Pi C\to B\Pi C'$（$\hat\varphi\colon E\Pi C\to E\Pi C'$）为 $A_{\infty}$（$L_{\infty}$）写像。由 $\varphi_*(b)=\sum_{k\ge1}\varphi_k(b^k)$（$\varphi_*(b)=\sum_{k\ge1}\frac{1}{k!}\varphi_k(b^k)$）定出 $\varphi_*\colon\hat M(C)\to\hat M(C')$。当 $\hat\varphi$ 是弱 $A_{\infty}$（$L_{\infty}$）写像时，也由 $\varphi_*(b)=\sum_{k\ge0}\varphi_k(b^k)$（$\varphi_*(b)=\sum_{k\ge0}\frac{1}{k!}\varphi_k(b^k)$）定出 $\varphi_*\colon\hat M(C)\to\hat M(C')$。
\end{definition}

\begin{definition}{6.132}
  当 $A_{\infty}$（$L_{\infty}$）代数与满足 $m_k\equiv0$（$k=0,1,2,\dots$）的 $A_{\infty}$（$L_{\infty}$）代数同伦同值时，称之为\textbf{形式的}（formal）。当弱 $A_{\infty}$（$L_{\infty}$）代数满足 $M(C)\neq\emptyset$ 时，称之为\textbf{非障碍的}（unobstructed）。

% [原文 p.308]
\end{definition}

在下面的定义中设 $\Lambda=\C[[T]]\otimes\C[u,u^{-1}]$。设 $C$ 是 $\Lambda$ 系数的滤过 $A_{\infty}$、$L_{\infty}$ 代数，$b\in\hat M(C)$。置 $b=\sum_{k\ge1}b_kT^k$。看 $\hat de^b$ 的 $T$ 的系数便知 $m_1(b_1)=0$。

\begin{definition}{6.133}
  称映射 $b\mapsto[b_1]$，即 $\hat M(C)\to H^0(\Pi\overline C\otimes\C[u,u^{-1}],m_1)$，为\textbf{小平–Spencer 写像}（Kodaira–Spencer map）。
\end{definition}

\begin{remark}{6.134}
  由规范等价的定义也容易看出，$b\mapsto[b_1]$ 诱导映射 $M(C)\to H^0(\Pi\overline C\otimes\C[u,u^{-1}],m_1)$。
\end{remark}

\begin{remark}{6.135}
  设 $\mathcal{M}(M)$ 是紧流形 $M$ 上复结构的模空间，$[\bar\partial]\in\mathcal{M}(M)$，用 (6.34) 的记号考虑 $\bar\partial+\epsilon b_1+\epsilon^2b_2+\cdots\mapsto[b_1]\in H^{1}_{\bar\partial}(M;TM)$，便得到映射 $T_{[\bar\partial]}\mathcal{M}(M)\to H^{1}_{\bar\partial}(M;TM)$。这才是本来意义的\textbf{小平–Spencer 写像}。
\end{remark}

\begin{theorem}{6.136}
  若 $C$ 是形式的，则小平–Spencer 写像 $\hat M(C)\to H_0(\Pi\overline C\otimes\C[u,u^{-1}],m_1)$ 是满射。
\end{theorem}

\begin{remark}{6.137}
  定理 6.107 是说，若给出障碍的群（在定理 6.136 的情况下是 $H^1(\Pi\overline C\otimes\C[u,u^{-1}],m_1)$）为 0，则小平–Spencer 写像是满射。定理 6.136 不假定这一条件。本节开头所做的计算表明，仓西写像可按 $H^0(\Pi\overline C\otimes\C[u,u^{-1}],m_1)$ 中元素的积与（高次）Massey 积展开为幂级数。$C$ 是形式的这一条件意味着积与（高次）Massey 积全为 0。因此在该条件下仓西写像消失。这就是定理 6.136。
\end{remark}

\begin{proof}[定理 6.136 的证明]
  设 $\hat\varphi\colon B\Pi C'\to B\Pi C$ 为同伦同值，且在 $B\Pi C'$ 上 $m_k$ 全为 0。此时对任意的
  \[
    [x]\in H^0(\Pi\overline C'\otimes\C[u,u^{-1}],\bar m_1)\simeq H^1(\overline C'\otimes\C[u,u^{-1}],m_1),
  \]
  置 $T[x]=b^0\in B\Pi C'$，则 $\hat de^{b^0}=0$。因此置 $b=\hat\varphi(e^{b^0})=\sum T^k\varphi_k(x^k)\in B\Pi C$，则 $\hat de^b=0$ 且 $[b_1]=[\varphi_1(x)]=[x]$。
\end{proof}

% CHUNK ch06b2 DONE

